Ternární model slouží k~vyhodnocení úspěšnosti koordinace zájmů mezi aktéry sociálního dialogu -- zaměstnanci, zaměstnavateli a~vládou. Každá z~těchto tří zájmových skupin je v~analytickém modelu popsána váženým indexem složeným z~dílčích statistických proměnných. Výsledná poloha země v~ternárním prostoru odráží relativní poměr a~rovnováhu sil mezi těmito třemi zájmovými směry. Váhy proměnných jsou stanoveny z~motivací vycházejících z~českého prostředí. Například v~zemích s~vyšším standardem mezd mohou být vedlejší faktory pro zaměstnance významnější.\par

\paragraph{Proměnné pro motivaci zaměstnanců}
zachycují vyjednávací úspěšnost a~podmínky zaměstnanců z~hlediska kupní síly, pracovní doby a~genderové rovnosti:
\begin{itemize}[noitemsep]
  \item \textbf{Medián disponibilního příjmu domácností (\SI{60}{\percent}):} medián ekvivalizovaného disponibilního příjmu domácností v~\acs{PPS} (zdroj Eurostat \texttt{ilc\_di01}, decil D5; viz obr.~\ref{fig:eu_prijem_distribuce}). Vyjadřuje, že kupní síla domácnosti představuje pro zaměstnance primární cíl kolektivního vyjednávání.
  \item \textbf{Skutečná týdenní odpracovaná doba (\SI{30}{\percent}):} do modelu je započtena inverzně (zdroj Eurostat \texttt{lfsa\_ewhan2}; viz tab.~\ref{tab:flexicurity}). Odráží fakt, že stejný příjem vykoupený delším časem v~práci představuje slabší vyjednaný výsledek.
  \item \textbf{Gender pay gap (\SI{10}{\percent}):} započten inverzně (zdroj Eurostat \texttt{earn\_gr\_gpgr2}; viz obr.~\ref{fig:problemy_gpg_sektor}). Slouží jako indikátor vnitřní spravedlnosti a~kvality odměňování.
\end{itemize}

\paragraph{Proměnné pro motivaci zaměstnavatelů}
odráží dostupnost pracovní síly, náklady na práci, daňové zatížení kapitálu a~provozní stabilitu z~hlediska firem:
\begin{itemize}[noitemsep]
  \item \textbf{Míra volných pracovních míst (\SI{40}{\percent}):} započtena inverzně (zdroj Eurostat \texttt{jvs\_a\_r21}; viz obr.~\ref{fig:stav_volna_mista_vyvoj}). Vyjadřuje napjatost trhu práce; vyšší podíl neobsazených míst zvyšuje náborové náklady a~představuje tlak na zaměstnavatele. Bez dostatku pracovní síly nemohou zaměstnavatelé efektivně konkurovat a~zvyšovat produkci.
  \item \textbf{Kombinovaná statutární sazba \aclgen{DPPO} (\SI{30}{\percent}):} započtena inverzně (zdroj \acs{OECD} \texttt{CTS\_CIT}; viz obr.~\ref{fig:stav_korporatni_dan_mapa}). Reprezentuje daňové zatížení výnosů.
  \item \textbf{Úplné hodinové náklady práce v~\acs{PPS} (\SI{20}{\percent}):} započteny inverzně (zdroj Eurostat \texttt{lc\_lci\_lev} a~\texttt{prc\_ppp\_ind}; viz tab.~\ref{tab:flexicurity}). Reprezentují cenovou mzdovou konkurenceschopnost práce.
  \item \textbf{Index sociálního míru odvozený ze stávkové aktivity (\SI{10}{\percent}):} hodnocený podle stávkových dnů (viz~část~\ref{par:spory}, obr.~\ref{fig:stav_kolektivni_spory_vyvoj}). Představuje korekci provozního rizika a~hrozby přerušení výroby.
\end{itemize}

\paragraph{Proměnné pro motivaci vlády}
zachycují politickou a~makroekonomickou stabilitu, fiskální kapacitu a~sociální výkonnost z~hlediska veřejné moci:
\begin{itemize}[noitemsep]
  \item \textbf{Kumulativní reálná změna \acs{HDP} 2020--2024 (\SI{40}{\percent}):} (zdroj Eurostat \texttt{nama\_10\_gdp}; viz obr.~\ref{fig:stav_hdp_vyvoj}). Odráží hospodářskou odolnost a~růstový potenciál země po krizových šocích, a~tedy i~politickou stabilitu.
  \item \textbf{\Acs{HDP} na obyvatele v~\acs{PPS} (\SI{30}{\percent}):} normováno vůči \acs{geo-EU27}=100 (zdroj Eurostat \texttt{nama\_10\_pc}; viz obr.~\ref{fig:stav_hdp_vyvoj}), což vyjadřuje celkový ekonomický výkon a~odráží výši daňových příjmů.
  \item \textbf{Míra zaměstnanosti ve věku 20–64 let (\SI{20}{\percent}):} (zdroj Eurostat \texttt{lfsi\_emp\_a}; viz obr.~\ref{fig:stav_zamestnanost}). Klíčová proměnná pro společenskou stabilitu a~financování sociálních systémů.
  \item \textbf{Giniho koeficient disponibilního příjmu (\SI{10}{\percent}):} započten inverzně (zdroj Eurostat \texttt{ilc\_di12}; viz obr.~\ref{fig:eu_gini_prijem}). Vysoká příjmová nerovnost snižuje sociální kohezi a~může ohrozit politickou stabilitu.\par
\end{itemize}

Výsledná vizualizace ternárního rozložení zemí a~příslušných silových polí (obr.~\ref{fig:practical_ternary_social_dialog} a~obr.~\ref{fig:practical_ternary_strength_map}) ilustruje odlišnost \acgen{ČR} od ostatních referenčních evropských modelů. V~diagramu se jasně formují tři odlišné shluky. Prvním je sociálně-demokratický model Dánska, který se spolu s~dalšími skandinávskými zeměmi se silným sociálním dialogem nachází nejblíže k~rovnováze. Vykazují vyváženou a~stabilní mzdově-koordinační strukturu balancující mezi produktivitou a~ochranou práce, jakkoli v~některých parametrech mírně inklinují k~pohledu konkurenceschopnosti a~zájmům zaměstnavatelů. Druhý shluk tvoří Německo a~Rakousko spolu s~Nizozemskem a~Belgií, kde mají zaměstnanci silnější pozici, což je dáno vysokým reálným disponibilním příjmem domácností při velmi nízké skutečné týdenní odpracované době a~silné státní podpoře zaměstnanců.\par

% Editable figure definition for practical_ternary_social_dialog
% Edit freely --- Python will NOT overwrite this file.
% To regenerate defaults, delete this file and re-run the script.
\def\strPracticalTernarySocialDialogTitle{Tripartitní rovnováha v~EU}%
\def\strPracticalTernarySocialDialogVertexA{Zam\v{e}stnanci}%
\def\strPracticalTernarySocialDialogVertexB{Firmy}%
\def\strPracticalTernarySocialDialogVertexC{Vl\'ada}%
%\def\strPracticalTernarySocialDialogCaption{Praktická tripartitní rovnováha v~\acs{EU} na základě modelu preferencí aktérů sociálního dialogu (odpovídající zaměření na zaměstnance, firmy a stát), 2026.}%
% --- Per-label y-nudge knobs (override with \renewcommand)
% Example: \renewcommand\NudgePracticalTernarySocialDialogCz{-3pt}  % shift label up by 3pt
\providecommand\NudgePracticalTernarySocialDialogCz{0pt}%
\providecommand\NudgePracticalTernarySocialDialogAt{0pt}%
\providecommand\NudgePracticalTernarySocialDialogDe{0pt}%
\providecommand\NudgePracticalTernarySocialDialogDk{0pt}%
\providecommand\NudgePracticalTernarySocialDialogPl{0pt}%
\providecommand\NudgePracticalTernarySocialDialogSk{0pt}%
\def\strPracticalTernarySocialDialogCaption{Ternární model rovnováhy naplnění potřeb aktérů sociálního dialogu, \acs{geo-EU27}, 2025. Model vážených veličin. Zdroj dat: Eurostat~\cite{eurostat_ilc_di01}~\cite{eurostat_lfsa_ewhan2_HR_weekly}~\cite{eurostat_gpg}~\cite{eurostat_jvs_a_r21}~\cite{eurostat_lc_lci_lev}~\cite{eurostat_prc_ppp_ind_PLI_GDP}~\cite{eurostat_nama_10_gdp}~\cite{eurostat_nama_10_pc}~\cite{eurostat_lfsi_emp_a}~\cite{eurostat_ilc_di12}; \acs{OECD}~\cite{oecd_cts_cit}; expertní skóre sociálního smíru~\cite{ilostat_STR_DAYS_ECO_RT_A}~\cite{CMKOS_ZpravaKV2025}.}%
%

% --- Auto-added angle nudge knobs in degrees (override with \renewcommand)
\providecommand\NudgeAnglePracticalTernarySocialDialogCz{25}%
\providecommand\NudgeAnglePracticalTernarySocialDialogDk{40}%
\providecommand\NudgeAnglePracticalTernarySocialDialogDe{22}%
\providecommand\NudgeAnglePracticalTernarySocialDialogSk{0}%
\providecommand\NudgeAnglePracticalTernarySocialDialogPl{8}%
\providecommand\NudgeAnglePracticalTernarySocialDialogAt{160}%
\begin{figure}[htb]
  \centering
  % Editable figure definition for practical_ternary_social_dialog
% Edit freely --- Python will NOT overwrite this file.
% To regenerate defaults, delete this file and re-run the script.
\def\strPracticalTernarySocialDialogTitle{Tripartitní rovnováha v~EU}%
\def\strPracticalTernarySocialDialogVertexA{Zam\v{e}stnanci}%
\def\strPracticalTernarySocialDialogVertexB{Firmy}%
\def\strPracticalTernarySocialDialogVertexC{Vl\'ada}%
%\def\strPracticalTernarySocialDialogCaption{Praktická tripartitní rovnováha v~\acs{EU} na základě modelu preferencí aktérů sociálního dialogu (odpovídající zaměření na zaměstnance, firmy a stát), 2026.}%
% --- Per-label y-nudge knobs (override with \renewcommand)
% Example: \renewcommand\NudgePracticalTernarySocialDialogCz{-3pt}  % shift label up by 3pt
\providecommand\NudgePracticalTernarySocialDialogCz{0pt}%
\providecommand\NudgePracticalTernarySocialDialogAt{0pt}%
\providecommand\NudgePracticalTernarySocialDialogDe{0pt}%
\providecommand\NudgePracticalTernarySocialDialogDk{0pt}%
\providecommand\NudgePracticalTernarySocialDialogPl{0pt}%
\providecommand\NudgePracticalTernarySocialDialogSk{0pt}%
\def\strPracticalTernarySocialDialogCaption{Ternární model rovnováhy naplnění potřeb aktérů sociálního dialogu, \acs{geo-EU27}, 2025. Model vážených veličin. Zdroj dat: Eurostat~\cite{eurostat_ilc_di01}~\cite{eurostat_lfsa_ewhan2_HR_weekly}~\cite{eurostat_gpg}~\cite{eurostat_jvs_a_r21}~\cite{eurostat_lc_lci_lev}~\cite{eurostat_prc_ppp_ind_PLI_GDP}~\cite{eurostat_nama_10_gdp}~\cite{eurostat_nama_10_pc}~\cite{eurostat_lfsi_emp_a}~\cite{eurostat_ilc_di12}; \acs{OECD}~\cite{oecd_cts_cit}; expertní skóre sociálního smíru~\cite{ilostat_STR_DAYS_ECO_RT_A}~\cite{CMKOS_ZpravaKV2025}.}%
%

% --- Auto-added angle nudge knobs in degrees (override with \renewcommand)
\providecommand\NudgeAnglePracticalTernarySocialDialogCz{25}%
\providecommand\NudgeAnglePracticalTernarySocialDialogDk{40}%
\providecommand\NudgeAnglePracticalTernarySocialDialogDe{22}%
\providecommand\NudgeAnglePracticalTernarySocialDialogSk{0}%
\providecommand\NudgeAnglePracticalTernarySocialDialogPl{8}%
\providecommand\NudgeAnglePracticalTernarySocialDialogAt{160}%
\begin{figure}[htb]
  \centering
  % Editable figure definition for practical_ternary_social_dialog
% Edit freely --- Python will NOT overwrite this file.
% To regenerate defaults, delete this file and re-run the script.
\def\strPracticalTernarySocialDialogTitle{Tripartitní rovnováha v~EU}%
\def\strPracticalTernarySocialDialogVertexA{Zam\v{e}stnanci}%
\def\strPracticalTernarySocialDialogVertexB{Firmy}%
\def\strPracticalTernarySocialDialogVertexC{Vl\'ada}%
%\def\strPracticalTernarySocialDialogCaption{Praktická tripartitní rovnováha v~\acs{EU} na základě modelu preferencí aktérů sociálního dialogu (odpovídající zaměření na zaměstnance, firmy a stát), 2026.}%
% --- Per-label y-nudge knobs (override with \renewcommand)
% Example: \renewcommand\NudgePracticalTernarySocialDialogCz{-3pt}  % shift label up by 3pt
\providecommand\NudgePracticalTernarySocialDialogCz{0pt}%
\providecommand\NudgePracticalTernarySocialDialogAt{0pt}%
\providecommand\NudgePracticalTernarySocialDialogDe{0pt}%
\providecommand\NudgePracticalTernarySocialDialogDk{0pt}%
\providecommand\NudgePracticalTernarySocialDialogPl{0pt}%
\providecommand\NudgePracticalTernarySocialDialogSk{0pt}%
\def\strPracticalTernarySocialDialogCaption{Ternární model rovnováhy naplnění potřeb aktérů sociálního dialogu, \acs{geo-EU27}, 2025. Model vážených veličin. Zdroj dat: Eurostat~\cite{eurostat_ilc_di01}~\cite{eurostat_lfsa_ewhan2_HR_weekly}~\cite{eurostat_gpg}~\cite{eurostat_jvs_a_r21}~\cite{eurostat_lc_lci_lev}~\cite{eurostat_prc_ppp_ind_PLI_GDP}~\cite{eurostat_nama_10_gdp}~\cite{eurostat_nama_10_pc}~\cite{eurostat_lfsi_emp_a}~\cite{eurostat_ilc_di12}; \acs{OECD}~\cite{oecd_cts_cit}; expertní skóre sociálního smíru~\cite{ilostat_STR_DAYS_ECO_RT_A}~\cite{CMKOS_ZpravaKV2025}.}%
%

% --- Auto-added angle nudge knobs in degrees (override with \renewcommand)
\providecommand\NudgeAnglePracticalTernarySocialDialogCz{25}%
\providecommand\NudgeAnglePracticalTernarySocialDialogDk{40}%
\providecommand\NudgeAnglePracticalTernarySocialDialogDe{22}%
\providecommand\NudgeAnglePracticalTernarySocialDialogSk{0}%
\providecommand\NudgeAnglePracticalTernarySocialDialogPl{8}%
\providecommand\NudgeAnglePracticalTernarySocialDialogAt{160}%
\begin{figure}[htb]
  \centering
  \input{../python/figures/practical_ternary_social_dialog.pgf}
  \caption{\strPracticalTernarySocialDialogCaption}
  \label{fig:practical_ternary_social_dialog}
\end{figure}

  \caption{\strPracticalTernarySocialDialogCaption}
  \label{fig:practical_ternary_social_dialog}
\end{figure}

  \caption{\strPracticalTernarySocialDialogCaption}
  \label{fig:practical_ternary_social_dialog}
\end{figure}


Třetím případem jsou postkomunistické země (Česko, Slovensko, Polsko), které vykazují společné strukturální zatížení, avšak s~odlišnými mzdově-koordinačními drahami. Zatímco Polsko je vzhledem k~většímu domácímu trhu a~odlišným odvětvovým tlakům taženo mírně odlišným směrem, Slovensko se vyznačuje polohou silně ovlivněnou extrémně nízkou příjmovou nerovností dle Giniho indexu, avšak za cenu nízkých reálných příjmů. Samotná \ac{ČR} je pak ze všech šesti zemí navzdory špatnému výsledku při stabilizaci ekonomiky v~obdobích krize nejvýrazněji vychýlena k~vládní ose makroekonomické stability, kumulativního růstu \acgen{HDP} a~extrémně vysoké zaměstnanosti.\par

Český sociální dialog funguje především jako formální a~institucionální kotva sociálního smíru a~vysoké participace, avšak při nízkých mzdových nákladech práce a~dlouhé pracovní době. Tento stav vyhovuje zájmům státu v~krátkodobé stabilizaci rozpočtů a~subdodavatelskému modelu nadnárodních korporací na úkor reálných příjmů tuzemských zaměstnanců. Jak ukazuje obr.~\ref{fig:practical_ternary_strength_map}, český systém sociálního dialogu vede v~průměru k~velmi nízkému uspokojení potřeb tripartitních aktérů.\par

% Editable figure definition for practical_ternary_strength_map
% Edit freely --- Python will NOT overwrite this file.
% To regenerate defaults, delete this file and re-run the script.
\def\strPracticalTernaryStrengthMapTitle{Souhrnné skóre sociálních partner\r{u}}%
\def\strPracticalTernaryStrengthMapColorbarLabel{skóre [1]}%
\def\strPracticalTernaryStrengthMapCaption{Průměr modelových os (A+B+C)/3 (souhrnné skóre sociálních partnerů), \acs{EU}. Zdroj dat: Eurostat.~\cite{eurostat_jvs_a_r21}~\cite{eurostat_lc_lci_lev}, \acs{OECD}.~\cite{oecd_cts_cit}}%
%
\begin{figure}[htbp]
  \centering
  \resizebox{\linewidth}{!}{% Editable figure definition for practical_ternary_strength_map
% Edit freely --- Python will NOT overwrite this file.
% To regenerate defaults, delete this file and re-run the script.
\def\strPracticalTernaryStrengthMapTitle{Souhrnné skóre sociálních partner\r{u}}%
\def\strPracticalTernaryStrengthMapColorbarLabel{skóre [1]}%
\def\strPracticalTernaryStrengthMapCaption{Průměr modelových os (A+B+C)/3 (souhrnné skóre sociálních partnerů), \acs{EU}. Zdroj dat: Eurostat.~\cite{eurostat_jvs_a_r21}~\cite{eurostat_lc_lci_lev}, \acs{OECD}.~\cite{oecd_cts_cit}}%
%
\begin{figure}[htbp]
  \centering
  \resizebox{\linewidth}{!}{% Editable figure definition for practical_ternary_strength_map
% Edit freely --- Python will NOT overwrite this file.
% To regenerate defaults, delete this file and re-run the script.
\def\strPracticalTernaryStrengthMapTitle{Souhrnné skóre sociálních partner\r{u}}%
\def\strPracticalTernaryStrengthMapColorbarLabel{skóre [1]}%
\def\strPracticalTernaryStrengthMapCaption{Průměr modelových os (A+B+C)/3 (souhrnné skóre sociálních partnerů), \acs{EU}. Zdroj dat: Eurostat.~\cite{eurostat_jvs_a_r21}~\cite{eurostat_lc_lci_lev}, \acs{OECD}.~\cite{oecd_cts_cit}}%
%
\begin{figure}[htbp]
  \centering
  \resizebox{\linewidth}{!}{\input{../python/figures/practical_ternary_strength_map.pgf}}
  \caption{\strPracticalTernaryStrengthMapCaption}
  \label{fig:practical_ternary_strength_map}
\end{figure}
}
  \caption{\strPracticalTernaryStrengthMapCaption}
  \label{fig:practical_ternary_strength_map}
\end{figure}
}
  \caption{\strPracticalTernaryStrengthMapCaption}
  \label{fig:practical_ternary_strength_map}
\end{figure}

